\chaptertitlepage{Literature Review}

This chapter reviews published work on deep learning for chest X-ray analysis, focusing on COVID-19, tuberculosis, and pneumonia, the three diseases this project detects. It covers the main CNN architectures, transfer learning, the datasets in common use, multi-disease systems, and the gap between accuracy in the lab and tools that clinicians can use. Where a study shaped a decision in our system, we say so. The architecture studies led us to DenseNet121 and EfficientNetB3 alongside a CNN trained from scratch, the transfer learning work led to our two-phase training, the augmentation and regularization studies set our shared training settings, and criticism of undeployed research led to the Flask ensemble application described in Chapters~3 and~4.

\section{Deep Learning in Medical Imaging}

Medical image analysis has moved from hand-crafted features and classical machine learning to end-to-end deep learning. Older pipelines needed experts to define features such as texture descriptors, histogram statistics, and edge detectors, which were then passed to a classifier such as a support vector machine or random forest. A CNN learns its features directly from the pixels, so no manual feature engineering is needed, and on well-defined screening tasks its performance is getting close to that of specialists. This matters for chest X-rays because disease often shows up as small changes in brightness and texture across the lungs. Fixed hand-crafted descriptors capture such patterns poorly, while convolutional filters learn them well.

\subsection{Foundational CNN Architectures}

Most pulmonary imaging research builds on a handful of CNN designs. Simonyan and Zisserman (2014) introduced the VGG networks and showed that stacking many small $3\times3$ filters extracts features well. Using several $3\times3$ layers in place of one larger kernel adds non-linearity and depth to the receptive field without too many parameters. VGG16 and VGG19 are still common baselines for binary chest X-ray classification, especially TB screening, where very high accuracy has been reported on public datasets. Their simple structure also makes them easy to fine-tune: the final fully connected layers are replaced with a new classification head.

He et al.\ (2016) proposed ResNet, whose skip connections solve the degradation problem, in which adding layers to a very deep network makes training error go up. Gradients flow directly through the identity paths, so networks with hundreds of layers train stably. ResNet50 and ResNet50V2 appear often in medical imaging benchmarks as general-purpose backbones that combine depth with stable training.

Huang et al.\ (2017) introduced DenseNet, in which each layer receives the feature maps of every earlier layer in its dense block, not only the layer just before it. This improves gradient flow, lets layers reuse features, and needs fewer parameters for a given depth. DenseNet works especially well on chest X-rays, where faint texture and opacity patterns need to survive through the network, which is why we chose DenseNet121 for the COVID-19 module.

Tan and Le (2019) proposed EfficientNet, which uses compound scaling to grow depth, width, and input resolution together by a fixed coefficient. Balancing all three, instead of enlarging one, gives good accuracy for the compute spent. EfficientNet suits deployments where inference time and memory matter, and those concerns led us to pick EfficientNetB3 for the tuberculosis module after benchmarking seven architectures.

\subsection{Large-Scale Chest X-Ray Benchmarks}

Deep learning for chest X-rays drew wide attention after CheXNet (Rajpurkar et al., 2017), a DenseNet121 model trained on ChestX-ray14 that detected pneumonia with radiologist-level accuracy on a held-out test set. CheXNet output a sigmoid probability of pneumonia and was trained on more than 100,000 frontal chest X-rays labelled with fourteen chest diseases. It made DenseNet a standard baseline for chest X-ray work and prompted a large body of follow-up research. It also showed that, on large datasets, deep learning can match specialists when the task is clearly defined and the evaluation is careful.

Irvin et al.\ (2019) released CheXpert, 224,316 chest X-rays labelled for 14 conditions by a rule-based labeler, with radiologists reviewing uncertain cases. Their DenseNet121 baseline matched radiologists on five conditions, including pneumonia and pleural effusion. It used an uncertainty-aware loss that handles ambiguous labels directly, instead of forcing them to positive or negative. CheXpert and similar large datasets showed that deep learning can be clinically useful for chest X-rays, provided that label noise, class imbalance, and differences between hospitals are handled through careful preprocessing and validation. CheXpert's size also shows that training one model on many conditions at once needs far more data than a student project has, which supports the modular single-task approach taken in this report.

These benchmark studies share three findings. First, almost all chest X-ray work starts from ImageNet weights, because medical datasets are rarely large enough to train deep networks from scratch without overfitting. Second, the choice of architecture affects both peak accuracy and inference cost, as our own TB benchmark showed when comparing VGG and EfficientNet. Third, the size and label quality of the dataset largely decide how well a model generalizes: high accuracy on a public Kaggle dataset does not carry over to images from another hospital or scanner without further validation. Our methodology follows from all three. We use transfer learning with phased fine-tuning for COVID-19 and TB, compare architectures systematically before choosing the TB model, and evaluate all three models on fixed held-out test splits.

\section{COVID-19 Detection from Chest X-Rays}

The COVID-19 pandemic sped up research into automated screening from portable chest X-rays. Emergency departments and poorly resourced hospitals needed fast triage when RT-PCR testing was scarce or slow. Portable X-ray units could be taken to ICU beds and field hospitals and produce images within minutes, so automated reading offered another route to a result when test kits or laboratories were overwhelmed.

Wang et al.\ (2020) introduced COVID-Net, a custom CNN that classifies X-rays as COVID-19, non-COVID pneumonia, or normal, and reported sensitivity above 90\%. They designed the architecture for the task, with depthwise separable convolutions and tailored filter banks, instead of relying only on standard ImageNet classifiers. Their results showed that COVID-19 opacities look similar to those of other pneumonias. That overlap remains a problem for multi-class screening and is one reason our system uses a single four-class label set.

Chowdhury et al.\ (2020) published the COVID-19 Radiography Database, one of the public sources merged into our manifest, and reported strong performance for COVID-19 versus normal on it. In our system these images supply the COVID-19 class, plus Normal and Pneumonia images where present. Images labelled ``Viral Pneumonia'' in the COVID sources go into the shared Pneumonia class, because the Kaggle pneumonia dataset mixes bacterial and viral cases and they cannot be separated again. The images are stored in one folder per class, so the manifest builder can load them by directory.

Rahman et al.\ (2021) studied how image enhancement and preprocessing affect COVID-19 detection and found that suitable contrast normalization and enhancement can raise performance by several percentage points. This supports using standard preprocessing matched to each model, which we do for all three models: \texttt{preprocess\_input} for DenseNet and EfficientNet, division by 255 for the custom CNN, and a resize to each model's expected input shape.

Zhang et al.\ (2020) introduced confidence-aware loss functions for COVID-19 detection that make predictions more reliable when labels are uncertain. Their model reached 96\% sensitivity while keeping false positives under control. Screening tools need both properties: during an outbreak a false negative can let the virus keep spreading, while too many false positives can fill the available isolation beds.

Taken together, the COVID-19 literature shows that DenseNet-family models, multi-phase fine-tuning, and careful separation of COVID-19 from other opacities are central to useful multi-class screening. We use DenseNet121, from the same family as CheXNet (Rajpurkar et al., 2017), as one of the three architectures trained on the four-class task. Our two-phase training, with a frozen backbone followed by partial unfreezing, follows the gradual fine-tuning recommended by Raghu et al.\ (2019). Making COVID-19 one of four shared output classes, instead of a separate single-disease app, addresses the comparability problem described in Section~\ref{sec:lit-gaps}. Chapter~3 describes the shared training code and Chapter~5 evaluates it.

\section{Tuberculosis Detection}

Tuberculosis is still one of the leading infectious causes of death. The World Health Organization (2022) reports more than 1.5 million TB deaths a year, and late diagnosis keeps the disease spreading, especially in high-burden regions with few radiologists. Chest X-rays are central to TB screening programmes, since upper-lobe infiltrates, cavities, and pleural effusion can suggest active disease before bacteriological tests confirm it. Automated TB detection from frontal chest X-rays is therefore a well-defined screening task with clear public health value.

Jaeger et al.\ (2014) set early TB detection benchmarks on the public Montgomery County and Shenzhen Hospital datasets, with standard evaluation splits that allowed algorithms to be compared reproducibly. Before these datasets, TB studies used inconsistent private data, so results across studies could not be compared reliably. The Montgomery and Shenzhen datasets are still reference points for TB screening research and influenced the later Kaggle TB datasets used in this project.

Lakhani and Sundaram (2017) applied an ensemble of GoogLeNet and AlexNet to chest X-ray datasets and reported AUC values of 0.99 and 0.97 respectively. Theirs was among the first studies to show that deep learning could match or beat radiologists at classifying TB from frontal chest X-rays. Their ensemble combined two architectures with different receptive fields, an early form of the multi-architecture comparison we carried out for the TB module.

Hwang et al.\ (2016) trained a CNN on more than 40,000 chest X-rays and reached an AUC of 0.96 for TB screening, before ImageNet transfer learning became common in medical imaging. Their work showed that with enough labelled examples, a larger dataset can make up for a simpler architecture.

More recent studies compare modern architectures on Kaggle TB datasets. Fine-tuned VGG models regularly reach very high accuracy on binary Normal versus Tuberculosis tasks, and EfficientNet models come close at a lower computational cost. Tan and Le (2019) showed that compound scaling lets EfficientNetB3 approach the accuracy of heavier architectures while running faster, which is why we included EfficientNetB3 in our four-class comparison alongside DenseNet121 and a custom CNN.

The TB literature supports two points we adopted. Architectures need to be benchmarked on the actual task, because the best model on an isolated binary task may not be the best on a harder shared multi-class problem. And binary TB screening on curated public datasets can reach near-perfect accuracy with transfer learning and regularization, but those figures cannot be compared with four-class results on patient-grouped splits. Real-world use would still need careful preprocessing, held-out evaluation, and prospective clinical validation.

\section{Pneumonia Detection}

CNN-based pneumonia detection has been benchmarked heavily on the Kaggle Chest X-Ray dataset released by Mooney (2018). It contains about 5,863 pediatric chest X-rays labelled NORMAL or PNEUMONIA, split into training, validation, and test folders. At that size overfitting is a known risk and regularization is needed. Because the splits are fixed, researchers can report results on the same held-out test set and compare architectures under identical conditions.

Studies on this dataset report accuracy above 90\% with architectures ranging from small sequential CNNs to large transfer learning models. Rajpurkar et al.\ (2017) and later CheXpert-related work treat pneumonia as one of the most studied chest X-ray findings, with CNNs learning the opacity, consolidation, and interstitial changes seen in bacterial and viral infection. On frontal pediatric X-rays, pneumonia usually appears as increased opacity in one or both lungs, often with air bronchograms inside the consolidated areas, and convolutional filters can learn these patterns when there are enough labelled examples.

On this moderately sized dataset, batch normalization, dropout, and learning rate scheduling consistently improve generalization. Batch normalization keeps activations stable across mini-batches, dropout switches off random neurons during training so they cannot co-adapt, and reducing the learning rate when validation loss plateaus allows finer weight updates. Custom CNNs of shallow to medium depth remain competitive when training time, model size, and a simple architecture matter more than the small accuracy gain of a very deep transfer model. That trade-off led to the five-block custom CNN used as Model~3.

Rahman et al.\ (2021) also studied preprocessing for pneumonia and COVID-19 classification and found that pixel normalization, resizing method, and color channel order all affect inference. Channel order matters when images are loaded with OpenCV and converted between BGR and RGB, as our shared preprocessing code does for all models. When preprocessing at inference differs from preprocessing in training, deployed medical AI systems can lose accuracy without any visible error.

The Kaggle pneumonia dataset labels images only as NORMAL or PNEUMONIA, with no subtype or cause. We map those folders into the shared Normal and Pneumonia classes, alongside the TB and COVID sources, so that every architecture is compared on the same four outputs.

\section{Transfer Learning}

When labelled medical data is limited, the standard approach is transfer learning from models pre-trained on ImageNet, such as VGG16, InceptionV3, ResNet50, DenseNet121, and EfficientNet. ImageNet has over a million natural images in a thousand categories, and a model trained on it learns general low-level features such as edges, textures, and shapes. These carry over well to X-rays, even though X-rays look very different from photographs. Medical datasets are rarely large enough to train deep networks from random weights without severe overfitting.

Raghu et al.\ (2019) studied transfer learning for medical imaging in their Transfusion paper. They found that low-level features transfer well, while higher-level features need fine-tuning for the new domain. They also observed that freezing the early layers and unfreezing deeper blocks gradually, or training a new classification head before fine-tuning the whole network, often beats training from random weights on small medical datasets. Keeping the backbone frozen at first stops large gradient updates from wrecking the pre-trained features while the new head is still random.

We follow this directly in the two-phase training of the DenseNet121 COVID-19 model. Phase~1 freezes the backbone and trains only the new head at a learning rate of $10^{-4}$. Phase~2 unfreezes the backbone from layer 300 onward and fine-tunes at $10^{-5}$. The EfficientNetB3 TB model works the same way: the ImageNet backbone starts frozen, and the classification head is trained with augmentation built into the model graph before the weights are saved for deployment.

The literature also warns that transfer learning does not guarantee generalization. Results on small public datasets can be inflated by leakage between training and test sets, by repeated tuning on the validation set, or by dataset artefacts such as hospital-specific borders or text that the model learns as a shortcut. Recommended practice is to fix random seeds, evaluate on a held-out test set, and compare architectures on identical splits, as we did in the TB benchmark with seed 1337.

\section{Data Augmentation and Regularization}

Medical imaging datasets are usually small compared with natural image datasets, so augmentation and regularization are needed to control overfitting. Augmentation enlarges the training set by applying transforms that do not change the label, which pushes the model to learn features that hold under those changes instead of memorizing pixels. For chest X-rays the literature favours random rotation and horizontal flips (frontal views are roughly symmetric left to right), width and height shifts that mimic small differences in patient position, zoom that mimics differences in scale, and other geometric transforms that leave intensities unchanged, applied only during training.

Dropout, batch normalization, early stopping, learning rate reduction on plateau, and checkpointing appear in most of the best-performing lung disease studies. Early stopping watches validation loss and stops training when it has not improved for a set number of epochs, so the model does not keep fitting noise. Checkpointing keeps the weights from the best epoch, since the final epoch may already be overfitted. Irvin et al.\ (2019) and Chowdhury et al.\ (2020) both pair augmentation with regularization when fine-tuning deep networks on imbalanced disease labels.

Class imbalance is common in pneumonia datasets, where positive cases can outnumber normal ones in some splits. For both sigmoid and softmax outputs it is usually handled by augmenting the minority class, splitting with stratification, and early stopping, instead of simply maximizing accuracy. Accuracy is misleading on imbalanced data: a model that always predicts the majority class can score well while missing every minority case. We apply these techniques to all three models, with augmentation settings adjusted to each model's input size and architecture. For example, DenseNet121 uses \texttt{ImageDataGenerator} with rotation and shift settings, while EfficientNetB3 has \texttt{RandomZoom} and \texttt{RandomRotation} layers built into the Keras model.

\section{Multi-Disease Detection Frameworks}

Far fewer studies build multi-disease systems than single-disease ones. There are many papers on COVID-19 detection, TB screening, or pneumonia classification, each reporting good accuracy on its own. Most of these tools handle one disease and report offline metrics, with no connection to a clinical workflow. A clinician who suspected several possible causes would need to run separate applications, each with its own interface, preprocessing, and output format. In emergency triage, where speed and simplicity count, that is impractical.

CheXpert and CheXNet predict many conditions at once from a single model trained on all of them. That approach needs very large labelled datasets, and it returns all predictions together instead of screening for a particular suspected condition. Clinical triage often works from a differential diagnosis: a clinician may suspect TB where it is endemic, COVID-19 during an outbreak, or community-acquired pneumonia in a child.

Our system addresses this by training three architectures (EfficientNetB3, DenseNet121, and a custom CNN) on one four-class dataset (Normal, Tuberculosis, COVID-19, Pneumonia), so that the comparison is fair and every model predicts the same diseases. Each backbone keeps its own input size and preprocessing (DenseNet $224\times224$, EfficientNet $256\times256$, custom CNN $150\times150$), but all share the same splits, seed, augmentation budget, and four-class softmax output, and the two transfer models share the same classification head. A Flask application runs every available model on each upload and combines their probabilities by soft voting. It shows where the models agree and disagree, and the user does not have to choose a disease model first. Published medical AI work rarely does this: most studies either report single-disease notebooks or train one large multi-label network on a dataset such as CheXpert, without a controlled comparison of backbones.

\section{Deployment and Clinical Translation}

A common criticism of medical AI research is the ``deployment gap'': many studies report accuracy, sensitivity, and AUC on retrospective datasets but never deliver software that clinicians can use. TensorFlow (Abadi et al., 2015) and its ecosystem made it practical to save models and run inference at scale, which lowered the engineering barrier. Even so, peer-reviewed chest X-ray studies often stop at a notebook, reporting test-set metrics without showing that a non-expert can run the model reliably through a stable interface. A model with 95\% accuracy in a controlled experiment gives no diagnostic value if it cannot be loaded, fed correctly preprocessed images, and run consistently in a clinic or community setting.

A deployable system needs engineering beyond training. It must check uploads in several stages and reject images that are not radiographs. Its inference preprocessing must match training exactly, including bilinear resizing where a library would default to nearest-neighbour. It should load large weight files only when needed so that the server starts quickly, and it must handle concurrent requests safely. OpenCV (Bradski, 2000) is the standard library for reading and checking images, and we use it in a three-stage chest X-ray gate (monochrome check, radiograph envelope features, and an anatomy veto) that runs before inference.

Lightweight Flask applications have become a common way for student and research prototypes to deploy medical AI. Our Flask application has an upload endpoint that accepts an image, validates it, runs \texttt{predict\_all()} over every trained four-class model, and shows each model's probability bars alongside a soft-voting ensemble, without asking the user to choose a disease model. Inference runs on the server, so model weights and TensorFlow stay there. Regulatory-grade clinical deployment is a separate matter from research prototypes, but a working interface still makes an academic project more credible.

Our Flask application, which runs all three four-class models and presents their ensemble, is our response to the deployment gap described in Section~\ref{sec:lit-gaps}.

\section{Explainability and Trust}

Clinicians will only adopt AI screening tools they can understand and trust, and accuracy alone does not earn that. A clinician given a prediction with no explanation is unlikely to act on it during triage, especially if it contradicts their own impression. Selvaraju et al.\ (2017) introduced Grad-CAM (Gradient-weighted Class Activation Mapping), which produces a heatmap of the image regions that most influenced a CNN's prediction. Grad-CAM takes the gradient of the target class score with respect to the activations of a chosen convolutional layer, giving a coarse localization map that can be laid over the X-ray.

Grad-CAM and similar saliency methods have been used on chest X-ray classifiers to check that models look at plausible anatomy, such as the lung fields, and not at dataset artefacts such as borders or text. If a model's heatmap lights up a hospital name printed on the border instead of the lung opacities, it has learned a shortcut that accuracy figures would never reveal. Despite wide interest, few multi-model screening systems with a single interface include explainability, and most explainability studies cover one model and one disease. Adding Grad-CAM or a similar visual explanation to a multi-disease system is an open opportunity, and we list it as future work in the gaps below and in Chapter~6.

\section{Gaps in Existing Research}
\label{sec:lit-gaps}

Deep learning for chest X-rays has advanced a great deal, but few results have become multi-disease screening tools that clinicians can use. Many studies report strong accuracy on curated public datasets, and few deliver integrated software that a clinician could use during triage. Each gap below shaped the design, architecture, and deployment work described in Chapters~3 to~5. Closing them takes better models, but also attention to software integration, validation discipline, and deployment, which the literature covers far less than new architectures.

\begin{enumerate}
    \item \textbf{Incomparable single-disease models:} Separate student/demo classifiers on incompatible label spaces prevent fair architecture comparison and fragment the user workflow.
    \item \textbf{Leakage-unaware splits:} Patient overlap and duplicate files across train and test inflate reported accuracy.
    \item \textbf{Accuracy-only reporting:} Headline accuracy hides weak minority-class recall (especially scarce COVID-19 samples).
    \item \textbf{Train--serve skew:} Different resize or normalisation at inference versus training silently degrades live predictions.
    \item \textbf{Weak input validation:} Monochrome-only checks accept many non-radiograph uploads.
    \item \textbf{Missing ensemble transparency:} Few demos surface multi-model disagreement alongside an aggregated verdict.
\end{enumerate}

\section{Summary}

The literature supports CNNs and ImageNet transfer learning for chest X-ray analysis. It also warns that small COVID datasets, dataset shortcuts, and reporting only accuracy weaken any claims made. Trustworthy results need clean data (patient-grouped splits and deduplication), matched training for a fair backbone comparison, and a deployable interface that states clearly it is not for diagnosis. Our project responds with a four-class benchmark across three architectures, evaluation that checks for leakage with a source probe, and a Flask ensemble application with a tested radiograph gate. Grad-CAM-style explainability is left for future work.

\begin{table}[H]
\centering
\caption{Selected related works versus this project's focus}
\label{tab:related-works}
\footnotesize
\begin{tabularx}{\textwidth}{@{}>{\RaggedRight\arraybackslash}p{2.2cm}>{\RaggedRight\arraybackslash}p{1.9cm}>{\RaggedRight\arraybackslash}p{2.3cm}>{\RaggedRight\arraybackslash}p{2.4cm}Y@{}}
\toprule
\HeaderRow \Head{Work} & \Head{Task} & \Head{Dataset style} & \Head{Method} & \Head{Limitation vs this FYP} \\
\midrule
CheXNet (2017) & Multi-label CXR & Large NIH corpus & DenseNet121 & Different label formulation; not a 4-class student demo \\
\ZebraRow COVID-Net (2020) & COVID multi-class & Custom CXR set & Custom CNN & Single architecture; pandemic-era bias risks \\
Kaggle pneumonia CNNs & Binary pneumonia & Mooney/RSNA-style & Shallow / transfer CNNs & Binary only; often leakage-unaware student splits \\
\ZebraRow Binary TB transfer studies & TB vs Normal & Public TB CXR & VGG / EfficientNet & Not comparable across diseases \\
\AccentRow \textbf{This project} & 4-class single-label & Merged public sources & 3 architectures + ensemble & Public small data; research demo only \\
\bottomrule
\end{tabularx}
\end{table}
